Medical foundation models are reused as frozen feature extractors, zero-shot classifiers and report
generators, so what their internal representations encode, and what the models actually use, matters
independently of any downstream head. This survey reviews the work that measures or intervenes on those
representations in medical vision encoders, image--text models and multimodal language models. It
organises the methods by family and by the point along the model at which they read the representation,
and it separates two questions that the literature often runs together: whether a concept can be decoded
from a representation, and whether the model uses it. Across \Ncore{} studies, clinical findings,
anatomy, demographic attributes and acquisition conditions are decodable from frozen encoders in four to seven modalities, but controls on probe validity are rare and no study shows that a model's output
depends on a named concept. The site and scanner an image comes from is the property recovered most
consistently, and whether removing it improves generalisation is untested. Comparisons of medical with
general pre-training never hold data, objective, architecture and scale fixed, so the effect of medical
pre-training remains unidentified. The survey closes with the experiments that would settle these
questions.
